\chapter{Dementia}

\section*{Definition}
Dementia is a syndrome of acquired, persistent decline in one or more cognitive abilities that interferes with independent daily functioning. Affected domains may include memory, attention, language, executive function, visuospatial skills, judgment, or social behavior. Dementia is not a normal part of aging and has many possible causes.

\section*{What Happens in Your Body}
Dementia results from disease or injury affecting brain networks. Neurodegenerative disorders may cause abnormal protein accumulation and progressive neuronal loss. Vascular disease can damage brain tissue through strokes or chronic small-vessel injury. Other conditions can impair cognition through inflammation, infection, metabolic disturbance, medication effects, sensory loss, or pressure within the skull. Different causes produce different patterns of cognitive, behavioral, motor, and functional change.

\section*{Causes}
\begin{itemize}
  \item Alzheimer disease, the most common cause of progressive dementia
  \item Vascular cognitive impairment or vascular dementia after strokes or small-vessel brain disease
  \item Dementia with Lewy bodies, often associated with visual hallucinations, fluctuating alertness, REM sleep behavior disorder, and parkinsonism
  \item Frontotemporal dementia, which may begin with personality, behavior, executive, or language changes rather than memory loss
  \item Parkinson disease dementia, Huntington disease, traumatic brain injury, and prion disease
  \item Potentially treatable contributors, including depression, medication effects, alcohol use, vitamin B12 deficiency, thyroid disease, sleep apnea, infection, hearing or vision loss, and metabolic disorders
\end{itemize}

\section*{When to See a Doctor}
Arrange medical assessment for progressive memory loss, repeated questions, getting lost, difficulty managing finances or medicines, impaired judgment, personality change, language difficulty, or loss of ability to perform familiar activities. Emergency care is required for sudden confusion, new weakness, facial droop, speech difficulty, seizure, severe headache, fever with neck stiffness, head injury, or a rapidly worsening change in mental status.

\section*{Related Symptoms}
\begin{itemize}
  \item Memory loss, disorientation, impaired attention, word-finding difficulty, or trouble planning and solving problems
  \item Personality or behavior change, apathy, agitation, depression, anxiety, delusions, or hallucinations
  \item Falls, gait change, tremor, rigidity, swallowing difficulty, incontinence, sleep disturbance, or reduced self-care
\end{itemize}

\section*{Self-Care / Home Management}
Maintain regular sleep, physical activity, social engagement, hearing and vision care, hydration, and a consistent daily routine. Use calendars, labels, pill organizers with supervision, medication lists, and fall-prevention measures. Do not leave a person at risk alone with driving, weapons, unsafe appliances, financial decisions, or medication administration. Caregivers should seek support and respite and should report abuse, neglect, wandering, or caregiver exhaustion.

\section*{How Your Doctor Will Evaluate You}
\begin{itemize}
  \item History should establish onset, progression, affected abilities, functional change, mood, sleep, substance use, medications, medical illness, family history, and collateral information from someone who knows the patient.
  \item Examination includes neurologic, mental-status, gait, sensory, hearing, vision, nutritional, and medication assessment.
  \item Cognitive screening may use the Mini-Cog, Montreal Cognitive Assessment, or Mini-Mental State Examination, but screening alone does not establish the cause.
  \item Laboratory evaluation commonly includes blood count, metabolic and liver testing, thyroid function, vitamin B12, and targeted tests for infection or other suspected causes.
  \item Brain MRI or CT evaluates vascular injury, tumor, hydrocephalus, subdural hematoma, and other structural disease. Specialist assessment, neuropsychological testing, genetic testing, cerebrospinal-fluid biomarkers, or amyloid imaging may be appropriate in selected cases.
\end{itemize}

\section*{Treatment Options}
Treatment begins with identifying and correcting reversible contributors, reviewing medications, treating depression or sleep apnea, correcting sensory loss, and managing vascular risk factors. Cholinesterase inhibitors may be used for Alzheimer disease, dementia with Lewy bodies, and Parkinson disease dementia; memantine may be used in selected moderate-to-severe Alzheimer disease. Behavioral and psychological symptoms are managed first with environmental and non-drug approaches. Antipsychotics require specialist-level caution because some people with dementia with Lewy bodies have severe sensitivity and all antipsychotics can increase serious adverse outcomes in older adults with dementia.

Selected patients with early Alzheimer disease may be evaluated for disease-modifying anti-amyloid therapy. These treatments require specialist confirmation of disease biology, MRI and bleeding-risk assessment, and monitoring for amyloid-related imaging abnormalities. Advance care planning, assessment of decision-making capacity, occupational support, caregiver education, and individualized safety planning are essential components of care.

\section*{References}
\begin{thebibliography}{99}
\bibitem{dementia:ref1} Livingston G, et al. Dementia prevention, intervention, and care: 2024 report of the Lancet standing Commission. \textit{Lancet}. 2024;404:572--628.
\bibitem{dementia:ref2} National Institute for Health and Care Excellence. Dementia: assessment, management and support for people living with dementia and their carers. NICE guideline NG97, updated 2025.
\bibitem{dementia:ref3} McKhann GM, et al. The diagnosis of dementia due to Alzheimer disease: recommendations from the National Institute on Aging and the Alzheimer's Association. \textit{Alzheimers Dement}. 2011;7:263--269.
\bibitem{dementia:ref4} American Psychiatric Association. \textit{The Diagnostic and Statistical Manual of Mental Disorders}. 5th ed., text revision. American Psychiatric Association; 2022.
\bibitem{dementia:ref5} U.S. Food and Drug Administration. Prescribing information and safety communications for disease-modifying therapies for Alzheimer disease. Accessed 2026.
\end{thebibliography}
